\documentclass[11pt]{article}

\usepackage[margin=1in]{geometry}
\usepackage{amsmath,amssymb,amsthm,mathtools}
\usepackage{booktabs}
\usepackage{graphicx}
\usepackage{subcaption}
\usepackage{placeins}
\usepackage{microtype}
\usepackage{natbib}
\usepackage[hidelinks]{hyperref}
\usepackage[nameinlink,noabbrev]{cleveref}

\newtheorem{theorem}{Theorem}[section]
\newtheorem{proposition}[theorem]{Proposition}
\newtheorem{lemma}[theorem]{Lemma}
\newtheorem{corollary}[theorem]{Corollary}
\theoremstyle{definition}
\newtheorem{definition}[theorem]{Definition}
\theoremstyle{remark}
\newtheorem{remark}[theorem]{Remark}

\newcommand{\R}{\mathbb{R}}
\newcommand{\E}{\mathbb{E}}
\newcommand{\cL}{\mathcal{L}}
\newcommand{\cF}{\mathcal{F}}
\newcommand{\cP}{\mathcal{P}}
\newcommand{\norm}[1]{\left\lVert #1\right\rVert}
\newcommand{\ip}[2]{\left\langle #1,#2\right\rangle}
\newcommand{\od}{\mathbin{\odot}}
\newcommand{\dd}{\,\mathrm{d}}
\newcommand{\rank}{\operatorname{rank}}
\newcommand{\diag}{\operatorname{diag}}

\title{How Much Can We Compress Neural Training Dynamics?\\
\large Response-memory closures for deep nonlinear feature learning}
\author{Draft research report}
\date{September 2026}

\begin{document}
\maketitle

\begin{abstract}
Compressing a trained function is easier than compressing its training dynamics:
two parameter states can define the same input--output map yet have different
tangent features and separate as soon as training resumes.  We study a more
demanding target: an autonomous state whose size does not grow with elapsed
training time and whose trajectory approximates gradient flow on every prescribed
finite time interval.  The setting keeps fixed depth, nonlinear activations and
hidden-feature learning under the canonical width-dependent mobilities.

The exact learned part of every hidden matrix is an integral of a backward
response against a forward response.  Exact history representations retain all
past responses and therefore do not compress this integral in time.  Our closure
stores only the first $P$ shifted-Legendre moments of each response history, while
retaining the initialized matrices exactly.  Moment differentiation closes because
the derivative of a Legendre mode is triangular in lower modes.  An ordinary
residual-activity clock gives a fully autonomous closure with compact-horizon
trajectory error $O_T(P^{-1})$.  A response-speed clock, which parametrizes the
joint forward/backward response path by residual mass plus path length, gives
$O_T(P^{-2})$ under one additional activation derivative.  Both statements hold
at fixed finite width and data for arbitrary fixed depth; they require neither
successful fitting nor special input geometry.  The second rate is an asymptotic
guarantee, not a claim of lower practical cost at small order.

The moving learned state falls from $(H-1)n^2$ dense entries to
$2(H-1)mnP$ history entries, but every initialized $n\times n$ matrix and its
forward/transpose actions remain.  Existing circle and MNIST experiments show
that very small orders can reproduce the dense fitted predictor, including on
passive inputs.  The report also separates this result from a width-uniform
population theorem, unrestricted scalar nonclosure, all-time uniform tracking,
and full compression of the initialized random operator, none of which is proved
here.
\end{abstract}

\section{The question is dynamical, not only functional}

The primary question is simple to state:
\emph{how much can the training dynamics of a neural network be compressed?}
A snapshot compressor need only reproduce the current function.  A training
compressor must also reproduce how the function will move.  Equality of
predictions does not identify the parameter-space Jacobian, the backward
responses, or the optimizer metric.  Consequently, two exact functional copies
can acquire different velocities and their errors can accumulate after training
continues.

This stricter question can be scientifically useful.  A successful compression
identifies a smaller state that retains the mechanism responsible for learning.
It may therefore offer a better set of coordinates for studying deep learning,
rather than only a smaller file format for a trained predictor.

We impose three non-negotiable features: more than one hidden layer, a
nonlinearity that remains active during the limiting description, and motion of
hidden representations at leading order.  The canonical mobilities used below
are of the same feature-learning type emphasized in tensor-program and mean-field
work \citep{yang2021tensor,nguyen2023rigorous}.  Population descriptions are
attractive because width becomes a particle approximation and many macroscopic
observables can become deterministic.  They do not automatically become
finite-dimensional: reused random matrices can survive as operators, and exact
descriptions may retain entire one- or two-time response histories.  Such a
description organizes the limit, but it does not by itself provide state whose
size is constant in training time.  Particle and measure gradient flows provide
an important shallow precedent \citep{chizat2018global}.

The motivating analogy is with statistical physics: more particles can permit a
smoother field description even while the microscopic system grows.  A more
speculative analogy is with models of computation.  A well-posed population flow
might play the role of an idealized model of learning, while finite networks
approximate it with finite width, finite precision and finite time steps.  These
analogies motivate the problem; they are not mathematical claims of this report.

\subsection{Claim levels and state classes}

We use the following distinctions throughout.

\begin{definition}[Compact-horizon global approximation]
Let $z(t)$ be a reference flow and $z_P(t)$ an autonomous order-$P$ closure.
The closure is \emph{globally accurate on compact horizons} if, for every fixed
$T<\infty$,
\[
  \sup_{0\le t\le T} d(z_P(t),z(t))\longrightarrow0
  \qquad(P\to\infty).
\]
The order may depend on $T$.  This does not mean that one fixed order works for
all $t\ge0$, nor that the constants remain bounded as $T\to\infty$.
\end{definition}

Three kinds of autonomous state are relevant.
\begin{enumerate}
\item A \emph{scalar ODE} evolves finitely many real numbers.  Its dimension is
independent of width, training time and any discretization of a population.
\item A \emph{population or field ODE} evolves finitely many functions, random
variables or densities on fixed source spaces.  It has finitely many fields but
infinitely many scalar degrees of freedom.
\item An \emph{operator ODE} evolves maps between population spaces.  A single
operator retains the row--column incidence that a finite collection of marginal
populations may lose.  ``One operator'' therefore need not mean low information.
\end{enumerate}
Exact and approximate closure are separate questions in each class.  So are a
local velocity match and uniform trajectory control on $[0,T]$.

Two established comparison results help locate the present contribution.  First,
for a one-input, three-hidden-layer linear feature-learning system, companion
theory constructs a fixed Hilbert space, a bounded source $\mathcal C_0$ and a
unique global trace-class increment $Q$ satisfying
\[
 \mathcal C=\mathcal C_0+Q,\qquad
 f=\tfrac14\operatorname{Tr}\mathcal C^4,\qquad
 \dot Q=-2\kappa(f-y)(\mathcal C_0^*+Q^*)^3.
\]
Finite gradient flow, and raw gradient descent for every step sequence tending to
zero, converge to its observables on compact horizons.  This is a scoped example
of exact operator closure, not a general deep nonlinear theorem.

Second, in the same benchmark there is no state-universal, width-uniform encoder
built from a fixed finite alphabet of bounded-degree contraction graphs that has
a polynomial autonomous scalar ODE (or a finite-order local polynomial PDE) and
a polynomial readout.  The proof differentiates the output repeatedly: the
$(2j-1)$st derivative contains a connected contraction
\[
 u^\top(B^\top R^\top RB)^j u
\]
of degree $4j+2$.  A fixed bounded-degree alphabet cannot generate these
pairwise distinct connected graphs.  This is evidence against simple exact
scalar closures.  It is not an impossibility theorem for arbitrary encoders;
indeed an unrestricted field can encode a whole future output profile.  We
therefore treat broad scalar impossibility and exact operator-to-population
compression as open questions rather than conclusions.

\section{Dense deep gradient flow}
\label{sec:dense}

Fix a dataset $(x_a,y_a)_{a=1}^m$, $x_a\in\R^d$, a width $n$, and $H\ge2$
hidden layers.  Set $v_a=x_a/\sqrt d$.  For simplicity every hidden layer has
width $n$; activations may differ by layer:
\begin{align}
 z_a^{(1)}&=W^{(1)}v_a,& h_a^{(1)}&=\phi_1(z_a^{(1)}),\nonumber\\
 z_a^{(\ell)}&=W^{(\ell)}h_a^{(\ell-1)},&
 h_a^{(\ell)}&=\phi_\ell(z_a^{(\ell)}),\quad 2\le\ell\le H,\\
 f_a&=\frac{w^\top h_a^{(H)}}{n},& r_a&=f_a-y_a,\qquad
 \rho=\left(\frac1m\sum_a r_a^2\right)^{1/2}.\nonumber
\end{align}
The loss is $\cL=\rho^2$.  Backward responses exclude the residual:
\begin{align}
 \delta_a^{(H)}&=w\od\phi_H'(z_a^{(H)}),\nonumber\\
 \delta_a^{(\ell)}&=\phi_\ell'(z_a^{(\ell)})\od
 (W^{(\ell+1)})^\top\delta_a^{(\ell+1)}.
\end{align}
With mobilities $(n,1,\ldots,1,n)$, the physical gradient field $F$ is
\begin{align}
 F_1&=-\frac2m\sum_a r_a\delta_a^{(1)}v_a^\top,\nonumber\\
 F_\ell&=-\frac{2}{nm}\sum_a r_a\delta_a^{(\ell)}
 (h_a^{(\ell-1)})^\top,\quad 2\le\ell\le H,\label{eq:dense-flow}\\
 F_w&=-\frac2m\sum_a r_a h_a^{(H)}.\nonumber
\end{align}
Thus every internal matrix has the exact representation
\begin{equation}
 W^{(\ell)}(t)=W_0^{(\ell)}-\frac{2}{nm}\sum_a
 \int_0^t r_a(s)\delta_a^{(\ell)}(s)
 (h_a^{(\ell-1)}(s))^\top\dd s.                 \label{eq:exact-history}
\end{equation}
The growing history in \eqref{eq:exact-history} is the object we compress.
The initialized matrices $W_0^{(\ell)}$ are retained exactly.

\section{Variant A: the residual-activity clock}
\label{sec:old-clock}

Let $p_k$ be shifted Legendre polynomials on $[0,1]$, normalized by
$p_k(1)=1$ and
\[
 \int_0^1p_j(x)p_k(x)\dd x=\frac{\mathbf 1_{j=k}}{2k+1}.
\]
They satisfy the triangular identity
\begin{equation}
 x p_k'(x)=k p_k(x)+\sum_{j<k}(2j+1)p_j(x).          \label{eq:triangular}
\end{equation}

Define the activity coordinate
\begin{equation}
 \tau(t)=1+\int_0^t\rho(s)\dd s.                    \label{eq:old-clock}
\end{equation}
At internal link $\ell$ and sample $a$, define
\[
 a_{\ell a}=h_a^{(\ell-1)},\qquad
 b_{\ell a}=\frac{r_a}{\rho}\delta_a^{(\ell)}\quad(\rho>0).
\]
In the coordinate $\xi=\tau(s)$, extend $a_{\ell a}$ constantly over the
artificial prefix $[0,1]$ and set $b_{\ell a}=0$ there.  Store only
\begin{align}
 H_{\ell ak}&=\int_0^\tau \bar a_{\ell a}(\xi)
 p_k(\xi/\tau)\dd\xi,&
 U_{\ell ak}&=\int_0^\tau \bar b_{\ell a}(\xi)
 p_k(\xi/\tau)\dd\xi,\quad 0\le k<P.               \label{eq:old-moments}
\end{align}
Differentiating and using \eqref{eq:triangular} gives the closed ODE
\begin{align}
 \dot H_{\ell ak}&=\rho a_{\ell a}-\frac{\rho}{\tau}
 \left[kH_{\ell ak}+\sum_{j<k}(2j+1)H_{\ell aj}\right],\nonumber\\
 \dot U_{\ell ak}&=r_a\delta_a^{(\ell)}-\frac{\rho}{\tau}
 \left[kU_{\ell ak}+\sum_{j<k}(2j+1)U_{\ell aj}\right]. \label{eq:old-ode}
\end{align}
Initially $H_{\ell a0}=a_{\ell a}(0)$ and all other moments vanish.  The
reconstructed matrices are
\begin{equation}
 \widehat W^{(\ell)}=W_0^{(\ell)}-\frac{2}{nm\tau}
 \sum_{a=1}^m\sum_{k=0}^{P-1}(2k+1)
 U_{\ell ak}H_{\ell ak}^\top.                       \label{eq:old-recon}
\end{equation}
All responses in \eqref{eq:old-ode} are evaluated in the current reconstructed
network.  The first matrix and readout evolve by their exact blocks in
\eqref{eq:dense-flow}.  Equations \eqref{eq:old-clock}--\eqref{eq:old-recon}
are therefore an autonomous finite ODE.  They never divide by $\rho$; the
normalized $b$ is only a proof coordinate.

If $\Pi_P$ denotes componentwise $L^2(0,\tau)$ projection onto polynomials of
degree less than $P$, then \eqref{eq:old-recon} is exactly
\begin{equation}
 \widehat W^{(\ell)}=W_0^{(\ell)}-\frac{2}{nm}\sum_a
 \int_0^\tau(\Pi_P\bar b_{\ell a})(\xi)
 (\Pi_P\bar a_{\ell a})(\xi)^\top\dd\xi.            \label{eq:projection-form}
\end{equation}
Thus closure of the moment ODE is exact; truncating the two history projections
is the modeling approximation.

\subsection{Compact-horizon convergence}

Let $\theta=(W^{(1)},\ldots,W^{(H)},w)$ and use the Euclidean block norm, with
Frobenius norms on matrices.

\begin{theorem}[Ordinary-clock response-memory convergence]
\label{thm:old}
Fix finite $n,d,m,H$, finite data and finite initialization.  Suppose
$\phi_\ell\in C^{1,1}_{\mathrm{loc}}(\R)$ for every layer.  For every
$T<\infty$ there are finite constants $C_{\mathrm{old}}(T)$ and $P_0(T)$,
independent of $P$, such that the order-$P$ closure exists uniquely through
$T$ for every $P\ge P_0(T)$ and
\begin{equation}
 \sup_{0\le t\le T}\norm{\widehat\theta_P(t)-\theta(t)}
 \le \frac{C_{\mathrm{old}}(T)}{\sqrt{P(P+1)}}.       \label{eq:old-rate}
\end{equation}
The same rate holds for all hidden responses and predictions on the training
set and on any fixed bounded collection of passive inputs.

If, in addition, every $\phi_\ell$ and $\phi_\ell'$ is globally bounded, then
the closure exists uniquely for all finite physical times for every $P\ge1$,
and \eqref{eq:old-rate} holds for every $P\ge1$ after enlarging the constant.
This includes tanh, logistic sigmoid, arctangent and smooth bounded periodic
activations.  No fitting, Gaussian initialization, input separation or label
compatibility assumption is required.
\end{theorem}

The word ``global'' in this theorem means the whole interval $[0,T]$ for each
arbitrary finite $T$.  It does not give a $T$-independent constant.

\subsection{Proof of the ordinary-clock theorem}

We give the proof because it shows exactly where history regularity enters and
why feedback does not invalidate the projection estimate.

\begin{lemma}[Legendre tail]
\label{lem:legendre}
For $q\in H^1(0,\tau;\R^N)$,
\begin{equation}
 \norm{q-\Pi_Pq}_{L^2(0,\tau)}^2
 \le\frac{1}{P(P+1)}\int_0^\tau \xi(\tau-\xi)
 \norm{q'(\xi)}^2\dd\xi
 \le\frac{\tau^2}{4P(P+1)}\norm{q'}_{L^2(0,\tau)}^2. \label{eq:leg-tail}
\end{equation}
\end{lemma}

\begin{proof}
On $[0,1]$, the shifted Legendre equation is
\[
 -\frac{\dd}{\dd x}\{x(1-x)p_k'(x)\}=k(k+1)p_k(x).
\]
Integration by parts shows that the derivatives are orthogonal in the weight
$x(1-x)$ and have squared norm $k(k+1)/(2k+1)$.  Bessel's inequality applied
to $q'$ in this weighted space gives
\[
 \sum_{k\ge1}k(k+1)\frac{\norm{c_k}^2}{2k+1}
 \le\int_0^1x(1-x)\norm{q'(x)}^2\dd x,
\]
where $c_k=(2k+1)\int_0^1q(x)p_k(x)\dd x$.  Parseval, polynomial density and
$k(k+1)\ge P(P+1)$ for $k\ge P$ give the first bound on $[0,1]$.
Rescaling and $\xi(\tau-\xi)\le\tau^2/4$ give
\eqref{eq:leg-tail}; summing componentwise proves the vector statement.
\end{proof}

\paragraph{Step 1: dense existence and a common compact region.}
The mobility matrix $\mathcal D$ has largest eigenvalue $n$, and
$F=-\mathcal D\nabla\cL$.  Hence
\begin{equation}
 \dot\cL=-\norm{\mathcal D^{-1/2}\dot\theta}^2,
 \qquad \int_0^t\norm{\dot\theta(s)}^2\dd s\le n\rho(0)^2. \label{eq:energy}
\end{equation}
Thus $\norm{\theta(t)-\theta(0)}\le\sqrt{nt}\rho(0)$.  A finite maximal
endpoint would have a finite limiting state, from which local existence extends
the solution.  Dense flow is therefore global.

Fix $T$ and let
\[
 R=\norm{\theta(0)}+\sqrt{nT}\rho(0)+1,
 \qquad \mathcal B=\{\vartheta:\norm{\vartheta}\le R\}.
\]
The dense path remains distance at least one from $\partial\mathcal B$.  Stop
the closure at its first exit from $\mathcal B$.  The rest of the proof uses
constants on this explicit ball and will rule out that exit.

For a concrete choice of constants set $X=\max_a\norm{v_a}$ and
$\alpha_0=X$.  Recursively choose
\begin{align}
 s_\ell&=\sup_{|z|\le R\alpha_{\ell-1}}|\phi_\ell'(z)|,&
 \alpha_\ell&=\sqrt n\sup_{|z|\le R\alpha_{\ell-1}}|\phi_\ell(z)|,\nonumber\\
 \beta_H&=R s_H,& \beta_\ell&=s_\ell R\beta_{\ell+1}\quad(\ell<H).
                                                               \label{eq:response-bounds}
\end{align}
Then $\norm{h_a^{(\ell)}}\le\alpha_\ell$ and
$\norm{\delta_a^{(\ell)}}\le\beta_\ell$ on $\mathcal B$.  Put
\begin{align}
 Y&=\left(m^{-1}\sum_a y_a^2\right)^{1/2},&
 q&=R\alpha_H/n+Y,& S&=Tq,& \mathsf A&=1+S,\nonumber\\
 v_1&=2\beta_1X,&
 v_\ell&=2\beta_\ell\alpha_{\ell-1}/n,& v_w&=2\alpha_H,\nonumber\\
 V&=\left(v_1^2+\sum_{\ell=2}^Hv_\ell^2+v_w^2\right)^{1/2}. \label{eq:compact-constants}
\end{align}
On the stopped interval, $\rho\le q$, $\tau\le\mathsf A$ and
$\norm{F}\le V\rho$.  Since $F$ is locally Lipschitz, it has some finite
Lipschitz constant $K$ on $\mathcal B$.

\paragraph{Step 2: an exact projection-error energy.}
For any one forward or backward history $q$ let
$D_q(t)=\norm{q-\Pi_Pq}_{L^2(0,\tau(t))}^2$ and let $q^*$ be the polynomial
projection evaluated at the current endpoint.  Differentiating the raw energy
and the projected energy, using \eqref{eq:old-ode}, cancels every basis-dilation
term and gives
\begin{equation}
 \dot D_q=\rho\norm{q-q^*}^2,
 \qquad D_q(t)=\int_0^t\rho(s)\norm{q(s)-q^*(s)}^2\dd s. \label{eq:projection-energy}
\end{equation}
The backward-prefix jump is harmless: only square integrability and growing
integrals are used.

Differentiate \eqref{eq:old-recon}.  The physical closure equation is exactly
\begin{equation}
 \dot{\widehat\theta}_P=F(\widehat\theta_P)+E,\qquad E_1=E_w=0,
\end{equation}
with internal defects
\begin{equation}
 E_\ell=\frac{2\rho}{nm}\sum_a
 (b_{\ell a}-b_{\ell a}^*)(a_{\ell a}-a_{\ell a}^*)^\top. \label{eq:defect}
\end{equation}
The two mixed projection terms vanish by orthogonality.  Consequently
\begin{equation}
 \int_0^t\norm{E_\ell(s)}_F\dd s
 \le\frac{2}{nm}\sum_a\sqrt{D_{b,\ell a}(t)D_{a,\ell a}(t)}. \label{eq:defect-energy}
\end{equation}
This is the bridge from a same-history approximation to a trajectory theorem.

\paragraph{Step 3: forward histories have a $P$-independent derivative
bound.}
Use primes for derivatives in $\xi=\tau(t)$ and define
\[
 Z_\ell(t)=\frac1m\sum_a\int_0^{\tau(t)}
 \norm{(h_a^{(\ell)})'(\xi)}^2\dd\xi.
\]
The first layer satisfies
$\norm{(h_a^{(1)})'}\le s_1v_1X$.  For later layers, the chain rule gives
\[
 (h_a^{(\ell)})'=\phi_\ell'(z_a^{(\ell)})\od
 \left[(F_\ell/\rho+E_\ell/\rho)h_a^{(\ell-1)}
 +\widehat W^{(\ell)}(h_a^{(\ell-1)})'\right].       \label{eq:forward-chain}
\]
Endpoint evaluation of a degree-$(P-1)$ polynomial has operator norm
$P/\sqrt\tau$ from $L^2(0,\tau)$, because
$\sum_{k<P}(2k+1)=P^2$.  Projection contraction and
$m^{-1}\sum_a\norm{b_{\ell a}}^2\le\beta_\ell^2$ imply
\[
 \left(m^{-1}\sum_a\norm{b_{\ell a}^*}^2\right)^{1/2}\le P\beta_\ell,
 \quad
 \left(m^{-1}\sum_a\norm{b_{\ell a}-b_{\ell a}^*}^2\right)^{1/2}
 \le(P+1)\beta_\ell.
\]
Combining this with \Cref{lem:legendre},
\eqref{eq:projection-energy} and \eqref{eq:defect} yields
\begin{equation}
 \int_0^t\rho\norm{E_\ell/\rho}_F^2\dd s
 \le\frac{2\beta_\ell^2\mathsf A^2}{n^2}
 Z_{\ell-1}(t).                                      \label{eq:e-square}
\end{equation}
The endpoint factor $P+1$ cancels the $P^{-1}$ projection factor.  This is
the step that prevents a loss of convergence order with depth.

Define finite, $P$-independent bounds recursively by
\begin{align}
 \overline Z_1&=S(s_1v_1X)^2,\nonumber\\
 \overline Z_\ell&=3s_\ell^2\left[
 S v_\ell^2\alpha_{\ell-1}^2+
 \left(R^2+\frac{2\beta_\ell^2\mathsf A^2\alpha_{\ell-1}^2}{n^2}\right)
 \overline Z_{\ell-1}\right].                       \label{eq:z-recursion}
\end{align}
Integrating \eqref{eq:forward-chain}, using
$\norm{x+y+z}^2\le3(\norm{x}^2+\norm{y}^2+\norm{z}^2)$ and
\eqref{eq:e-square}, proves inductively that $Z_\ell\le\overline Z_\ell$.

\paragraph{Step 4: accumulated defect and feedback.}
The zero backward prefix gives
$m^{-1}\sum_aD_{b,\ell a}\le S\beta_\ell^2$.  By
\Cref{lem:legendre},
\[
 \frac1m\sum_aD_{a,\ell a}
 \le\frac{\mathsf A^2\overline Z_{\ell-1}}{4P(P+1)}.
\]
Summing \eqref{eq:defect-energy} over links gives
\begin{equation}
 \int_0^t\norm{E(s)}\dd s
 \le\frac{B_{\mathrm{old}}}{\sqrt{P(P+1)}},\qquad
 B_{\mathrm{old}}=\frac{\mathsf A}{n}
 \sum_{\ell=2}^H\beta_\ell\sqrt{S\overline Z_{\ell-1}}. \label{eq:old-defect-bound}
\end{equation}
Subtracting dense flow from the perturbed flow and applying the Lipschitz
bound on $F$ gives
\[
 e(t)\le\frac{B_{\mathrm{old}}}{\sqrt{P(P+1)}}
 +K\int_0^te(s)\dd s.
\]
Gronwall's inequality therefore yields
\begin{equation}
 \sup_{t\le T}e(t)\le
 \frac{B_{\mathrm{old}}e^{KT}}{\sqrt{P(P+1)}}.       \label{eq:old-explicit}
\end{equation}
Choose $P_0$ so the right side is less than $1/2$.  The closure then cannot
reach $\partial\mathcal B$.  Its moment representations, $1\le\tau\le\mathsf A$
and the response bounds place its whole finite state in a compact subset of the
locally Lipschitz domain.  A finite maximal endpoint would therefore extend,
so the closure exists through $T$.  Lipschitz dependence of network responses
on $\theta$ gives the final prediction statement.

For completeness, suppose every activation and first derivative is globally
bounded.  Set
\[
 \alpha_\ell=\sqrt n\norm{\phi_\ell}_\infty,qquad
 s_\ell=\norm{\phi_\ell'}_\infty,qquad
 B_w=\sqrt{\norm{w(0)}^2+nY^2T}.
\]
The unchanged readout equation obeys
\[
 \frac{\dd}{\dd t}\norm w^2
 =nY^2-4n\frac1m\sum_a(f_a-y_a/2)^2\le nY^2,
\]
for both dense and closure flows.  Therefore
$\norm w\le B_w$, $\rho\le q:=B_w\alpha_H/n+Y$,
$S=Tq$ and $\tau\le\mathsf A:=1+S$.  Start with
$\beta_H=B_ws_H$ and proceed downward.  For $\ell=H,\ldots,2$, define
\[
 D_\ell=\norm{W_0^{(\ell)}}_F+
 \frac{2\mathsf A\beta_\ell\alpha_{\ell-1}}{n},
 \qquad \beta_{\ell-1}=s_{\ell-1}D_\ell\beta_\ell.
\]
Projection contraction in \eqref{eq:projection-form} gives
$\norm{\widehat W^{(\ell)}}_F\le D_\ell$; the backward recurrence then gives
$\norm{\delta^{(\ell-1)}}\le\beta_{\ell-1}$.  There is no circularity because
the descent starts from the bounded readout.  Finally
\[
 \norm{W^{(1)}(t)}_F\le\norm{W^{(1)}(0)}_F+2SX\beta_1.
\]
These $P$-independent physical bounds, the moment integral representations and
$1\le\tau\le\mathsf A$ put the order-$P$ state in a compact set for every
fixed $P$.  Hence the closure continues through $T$ for every $P\ge1$.
Repeating Steps 2--4 with these global constants gives
\eqref{eq:old-explicit} for every $P$.  This completes the proof.

\section{Variant B: a joint response-speed clock}
\label{sec:new-clock}

The activity clock controls the first-layer forward history particularly
cleanly, but it does not make every backward history smooth.  A more symmetric
coordinate parametrizes the entire encoded response path.  Concatenate
\begin{equation}
 \Psi(\theta)=\operatorname{stack}_{\ell=2,\ldots,H;\,a=1,\ldots,m}
 (a_{\ell a},b_{\ell a}),\qquad
 \dot L=g=\rho+\norm{\dot\Psi}_2,\quad L(0)=1.        \label{eq:new-clock}
\end{equation}
Then $\norm{\dd\Psi/\dd L}\le1$.  The clock changes where history is placed,
but the learning integral still carries residual mass.  At time $s$, place a
point at $\xi=L(s)$ and assign mass $\rho(s)\dd s$.  Give $[0,1]$ Lebesgue
mass and matching constant prefixes.  Thus
\[
 \int q\dd\mu_t=\int_0^1q(\xi)\dd\xi+
 \int_0^tq(L(s))\rho(s)\dd s,\qquad \dd\mu_t\le\dd\xi.
\]
The mass is $A(t)=1+\int_0^t\rho\dd s$, generally different from $L(t)$.

Let $p=(p_0,\ldots,p_{P-1})^\top$, $e=(1,\ldots,1)^\top$, and let
$\mathsf T_{kk}=k$, $\mathsf T_{kj}=2j+1$ for $j<k$.  Store one shared Gram
matrix and link/sample moments
\[
 H_{\ell a}=\int\bar a_{\ell a}p^\top\dd\mu,\qquad
 U_{\ell a}=\int\bar b_{\ell a}p^\top\dd\mu,\qquad
 G=\int pp^\top\dd\mu,
\]
where $p$ is evaluated at $\xi/L(t)$.  Put
$C_{\ell a}=b_{\ell a}(0)a_{\ell a}(0)^\top$.  The reconstruction and exact
moment equations are
\begin{align}
 \widehat W^{(\ell)}&=W_0^{(\ell)}-\frac{2}{nm}\sum_a
 (U_{\ell a}G^{-1}H_{\ell a}^\top-C_{\ell a}),       \label{eq:new-recon}\\
 \dot H_{\ell a}&=\rho a_{\ell a}e^\top-(g/L)H_{\ell a}\mathsf T^\top,
 &\dot U_{\ell a}&=r_a\delta_a^{(\ell)}e^\top-(g/L)U_{\ell a}\mathsf T^\top,
 \nonumber\\
 \dot G&=\rho ee^\top-(g/L)(\mathsf TG+G\mathsf T^\top). \label{eq:new-ode}
\end{align}
Initially $H_{\ell a}=[a_{\ell a}(0),0,\ldots]$,
$U_{\ell a}=[b_{\ell a}(0),0,\ldots]$ and
$G=\diag(1,(1/3),\ldots,(1/(2P-1)))$.  The prefix correction makes
$\widehat W^{(\ell)}(0)=W_0^{(\ell)}$.

Solve $Gq=e$ and set $a^*_{\ell a}=H_{\ell a}q$ and
$b^*_{\ell a}=U_{\ell a}q$.  Product differentiation cancels every $g/L$
term and gives the physical defect
\begin{equation}
 E_\ell=\frac{2\rho}{nm}\sum_a
 (b_{\ell a}-b^*_{\ell a})(a_{\ell a}-a^*_{\ell a})^\top. \label{eq:new-defect}
\end{equation}
This has the same product-of-errors form as Variant A, but now both factors
belong to one path with unit speed in $L$.

\begin{theorem}[Response-clock convergence]
\label{thm:new}
Under the finite setup of \Cref{thm:old}, assume
$\phi_1\in C^{1,1}_{\mathrm{loc}}$ and
$\phi_\ell\in C^{2,1}_{\mathrm{loc}}$ for $\ell\ge2$.  If $\rho(0)>0$, then
for every finite $T$ there are $P_0(T)$, $\Lambda_T$ and $C_T$, independent
of $P$, such that for $P\ge P_0(T)$ the ODE
\eqref{eq:new-clock}--\eqref{eq:new-ode} has a unique regular solution on
$[0,T]$, $L_P(t)\le\Lambda_T$, and
\begin{equation}
 \int_0^T\norm{E(t)}\dd t\le\frac{C_T}{P(P+1)},\qquad
 \sup_{t\le T}\norm{\widehat\theta_P(t)-\theta(t)}
 \le\frac{C_Te^{K_TT}}{P(P+1)}.                      \label{eq:new-rate}
\end{equation}
If $\rho(0)=0$, both flows are stationary.  The conclusion is
$O_T(P^{-2})$ at fixed width, depth and data.  It is not an all-small-order or
uniform-all-time statement.
\end{theorem}

\begin{proof}[Proof architecture]
For a history $q$, let $D_q$ be its weighted least-squares error.  Differentiating
$\operatorname{tr}(H G^{-1}H^\top)$ with \eqref{eq:new-ode} proves the exact
identity
\[
 \dot D_q=\rho\norm{q-q^*}^2.
\]
Consequently, Cauchy--Schwarz in physical time gives the analogue of
\eqref{eq:defect-energy}.  On a compact physical ball, raw history energies
give a coarse, order-independent bound on $\int\norm{E}$ and hence on total
physical path variation.  Dense residual satisfies a finite-horizon lower
bound $\rho_{\rm dense}\ge\mu_T>0$.  On the region
$\rho\ge\mu_T/2$, the response map has a finite derivative bound
$\norm{D\Psi}\le J_T$.  Therefore
\[
 L_P(T)\le A_T+J_T\left(\int_0^T\norm{F}\dd t+
 \int_0^T\norm{E}\dd t\right)=:\Lambda_T,
\]
without assuming a clock bound.

The joint extended history is $1$-Lipschitz in $\xi$ and constant on the
prefix.  Weighted best approximation is no worse than ordinary Legendre
projection because $\dd\mu\le\dd\xi$.  Applying \Cref{lem:legendre} to the
whole concatenated path yields
\[
 \sum_{\ell,a}(D_{a,\ell a}+D_{b,\ell a})
 \le\frac{L_P^2(L_P-1)}{4P(P+1)}.
\]
The arithmetic--geometric mean inequality in the defect estimate gives
\[
 \int_0^T\norm{E}\dd t
 \le\frac{\Lambda_T^2(\Lambda_T-1)}{4nmP(P+1)}.
\]
Gronwall gives \eqref{eq:new-rate}.  For sufficiently large $P$, this keeps
the closure inside the physical ball and keeps its residual above
$\mu_T/2$.  The prefix makes $G$ positive definite; bounded $L$ gives a
positive minimum eigenvalue for each fixed $P$.  All moment variables then
stay in a compact subset of the locally Lipschitz domain, which proves
continuation through $T$ and closes the first-exit argument.
\end{proof}

\subsection{The new clock is fully implementable}

Equation \eqref{eq:new-clock} appears implicit only because $\dot\Psi$ is
written compactly.  No nonlinear solve for $g$ is needed.  First compute
$q=G^{-1}e$, endpoint fits $a^*,b^*$, and internal velocities
\begin{equation}
 V_\ell=\dot{\widehat W}^{(\ell)}=-\frac{2\rho}{nm}\sum_a
 \left[b_{\ell a}(a^*_{\ell a})^\top+
 b^*_{\ell a}(a_{\ell a}-a^*_{\ell a})^\top\right]. \label{eq:internal-velocity}
\end{equation}
These expressions contain no $g$.  Compute $V_1=F_1$ and $V_w=F_w$, then a
directional forward pass
\[
 \dot z_a^{(1)}=V_1v_a,\qquad
 \dot z_a^{(\ell)}=V_\ell h_a^{(\ell-1)}+
 \widehat W^{(\ell)}\dot h_a^{(\ell-1)},\qquad
 \dot h_a^{(\ell)}=\phi_\ell'(z_a^{(\ell)})\od\dot z_a^{(\ell)}.
\]
This gives $\dot r$ and $\dot\rho$.  A reverse directional pass differentiates
the backward recurrence, using $\phi_\ell''$ for $\ell\ge2$, and then
\[
 \dot b_{\ell a}=\frac{\dot r_a\delta_a^{(\ell)}+
 r_a\dot\delta_a^{(\ell)}}{\rho}-b_{\ell a}\frac{\dot\rho}{\rho}.
\]
Now evaluate $g=\rho+\norm{\dot\Psi}$ and finally
\eqref{eq:new-ode}.  Matrix and transpose actions use the same reconstructed
operator.  One factors $G$ and solves linear systems; an explicit inverse is
unnecessary.

The response clock is more regular in approximation coordinates, but it is not
automatically better at low order.  At $P=1$ the polynomial space is constant,
so changing the coordinate alone leaves the physical constant-mode closure
unchanged.  Higher orders pay for Gram solves and derivative passes, and the
theoretical constants can be large.

\section{What is compressed?}

Let $H$ be hidden depth, $n$ width, $m$ training samples and $P$ moment order.
The comparison in \Cref{tab:complexity} separates moving learned state from
fixed source storage.

\begin{table}[t]
\centering
\caption{Leading storage and direct-operation counts.  A full batch applies
each internal map to $m$ vectors.  Counts omit lower-order activations and
first/readout layers.}
\label{tab:complexity}
\small
\begin{tabular}{lll}
\toprule
 & Dense gradient flow & Response-memory closure \\
\midrule
Moving internal state & $(H-1)n^2$ & $2(H-1)mnP$ \\
Fixed internal source & included above & $(H-1)n^2$ for all $W_0^{(\ell)}$ \\
Learned rank/link & up to $n$ & $\le mP$ (A), $\le m(P+1)$ (B) \\
One learned action & $O(n^2)$ & $O(nmP)$ \\
One full-batch action & $O(mn^2)$ & $O(nm^2P)$ \\
Moment dilation & --- & $O((H-1)mnP)$ \\
Variant-B overhead & --- & $O(P^3+(H-1)mnP^2)$ \\
\bottomrule
\end{tabular}
\end{table}

The learned action is matrix-free:
\[
 \widehat W^{(\ell)}v=W_0^{(\ell)}v-
 \frac{2}{nm}\sum_a U_{\ell a}G^{-1}H_{\ell a}^\top v
 \quad\text{(with the prefix correction in Variant B)},
\]
and the transpose uses the same factors in reverse order.  Storage is
independent of elapsed training steps.  The initialized operator remains
uncompressed and still costs $O(n^2)$ per vector under a direct dense action.
The current result is therefore compression of the \emph{learned moving
increment}, not total elimination of dense operators.  Total stored memory can
exceed a dense implementation if $W_0$ and all histories are materialized.

Ignoring constants and the order threshold, \Cref{thm:old} gives the sufficient
scaling $P=O(\varepsilon^{-1})$ for trajectory error $\varepsilon$, while
\Cref{thm:new} gives $P=O(\varepsilon^{-1/2})$.  Thus the proven moving-state
scalings are respectively
\[
 O(Hmn\varepsilon^{-1})\quad\text{and}\quad
 O(Hmn\varepsilon^{-1/2}).
\]
These are approximation guarantees, not practical speed laws: $C_T$ and
$P_0(T)$ depend on the task, width, depth and horizon, and Variant B has larger
per-step overhead.

\subsection{Population interpretation and its precise limit}

For fixed $P$, the neuron-indexed vectors can be viewed as particles sampling
population fields indexed by the initialization source.  This makes the closure
a natural Monte Carlo implementation of a hybrid population/operator system:
the response fields move, while the initialized random operator and its adjoint
remain exact.  The theorems above, however, keep $n$ fixed and their constants
may depend on $n$.

There is still a useful diagonal statement.  Suppose an independent theorem
establishes that dense width-$n$ flows $\theta_n$ converge on $[0,T]$ to a
population object $\Theta$ in a metric $d$, and suppose the observables entering
$d$ are controlled by \Cref{thm:old} or \Cref{thm:new}.  Choose any
$n_j\to\infty$, then choose $P_j$ large enough that the fixed-width closure
error is at most $1/j$.  The triangle inequality gives
\[
 d(\widehat\theta_{n_j,P_j},\Theta)
 \le d(\widehat\theta_{n_j,P_j},\theta_{n_j})+d(\theta_{n_j},\Theta)
 \longrightarrow0.
\]
This constructs a sequential approximation to the population flow.  It does
not prove that a fixed $P$ works uniformly in width, that the $W_0$ action has a
finite-population representation, or that the two limits commute.  Rigorous
multilayer mean-field and high-dimensional DMFT results exist in other regimes
\citep{nguyen2023rigorous,celentano2021highdimensional}; matching their
architectures, initializations and observables is a separate theorem.

\section{Why this is not a frozen dictionary}

A frozen dictionary chooses response directions at initialization and trains
only their coefficients.  It can be compact, but it cannot rotate its atoms to
follow newly learned features.  A trainable low-rank factorization allows more
motion, yet its factor-gradient geometry is still different from the integrated
rank-one updates of dense gradient flow.

Response memory keeps a low-rank learned correction at every instant, but its
atoms are moments of the \emph{current closure trajectory}.  The loop is
\[
 \text{moments}\longrightarrow\widehat W
 \longrightarrow\text{current forward/backward responses}
 \longrightarrow\dot{\text{moments}}.
\]
The basis in history coordinate is fixed; the neuron vectors multiplying that
basis coevolve nonlinearly.  This is why the method can represent substantial
feature learning without being a frozen tangent model.

On the hard two-outlier circle task at width $2048$, the current study reports
RMS discrepancies $1.100$ and $0.756$ for frozen derivative dictionaries with
18 and 72 stored vectors, and $0.896$ for a trainable 18-vector derivative
dictionary.  Response memory gives $0.07295$ at $P=3$ (48 vectors) and
$0.00466$ at $P=7$ (112 vectors).  In a more controlled comparison retaining
the same $W_0$, response moments beat a trained factorization $W_0+AB$ in all
29 valid task/order/seed cells; on the hard task at rank 56 the errors were
$0.004663$ versus $0.516210$ and $0.251026$.  These observations isolate a
mechanistic difference from those baselines.  They do not prove superiority
over every adaptive dictionary or at matched total memory.

\section{Empirical evidence}
\label{sec:experiments}

The experiments use finite networks and numerical time integration.  Circle
scores compare predictors on 8192 passive angles; MNIST scores use held-out
images.  Most comparisons match each model at its own first training-MSE
$0.001$ crossing.  They therefore test agreement of fitted functions, while
the theorems control same-time trajectories.  The panels are selected
finite-instance studies, generally with one initialization, and do not estimate
the asymptotic theorem constants.  The three figures below use the ordinary
activity clock.

Variant B has also been implemented with batched Gram solves and directional
response derivatives.  In a matched width-$2048$ comparison on three circle
tasks with tanh depth two and GELU depth three, all 49 dense/closure fits reached
training RMS at most $0.05$; its $P=3$ predictor discrepancy was below $0.05$ in
five of six principal groups.  The GELU alternating-quadrant exception improved
from $1.245$ to $0.112$ after halving the time step.  In a broader 82-case sweep,
the fitted-pair median whole-query discrepancies for $P=1,2,3$ were
$0.0427,0.0189,0.00680$, but several fitted alternating-label cases remained
poor.  These runs establish implementability and useful low-order cases.  They
do not empirically certify the $O(P^{-2})$ rate or show that Variant B uniformly
beats Variant A.

\begin{figure}[t]
 \centering
 \includegraphics[width=\textwidth]{figures/circle_shallow.pdf}
 \caption{Two-hidden-layer tanh networks of width $2048$.  The black curve is
 the dense fitted function on the whole circle; colored curves are ordinary-clock
 closures.  Training labels are marked by crosses.  Increasing order reproduces
 the dense predictor between and far away from training inputs.}
 \label{fig:circle-shallow}
\end{figure}
\FloatBarrier

Across the five tasks in \Cref{fig:circle-shallow}, the measured RMS errors are:
\begin{center}
\small
\begin{tabular}{lrrr}
\toprule
Task & $P=1$ & $P=3$ & $P=7$\\
\midrule
Two outliers, alternating & 0.2363 & 0.0730 & 0.00466\\
Quadrant, alternating & 0.3414 & 0.0457 & 0.00165\\
Quadrant, paired & 0.0193 & 0.00282 & $\approx0.00012$\\
Quadrant, center/edges & 0.0456 & 0.00487 & 0.00130\\
Equally spaced, mixed odd & 0.0124 & 0.000520 & $\approx0.00001$\\
\bottomrule
\end{tabular}
\end{center}
These tasks include rapid alternation and large unlabeled arcs, so
training-point agreement alone would not force the observed whole-circle match.

\begin{figure}[t]
 \centering
 \includegraphics[width=0.86\textwidth]{figures/circle_deep.pdf}
 \caption{Three-hidden-layer tanh networks of width $4096$, compared at matched
 training loss.  Across the five tasks, $P=3$ whole-circle RMS discrepancy ranges
 from $2.82\times10^{-4}$ to $5.85\times10^{-2}$.  The order trend is not
 pointwise monotone: $P=2$ beats $P=3$ on the two alternating tasks.}
 \label{fig:circle-deep}
\end{figure}
\FloatBarrier

The depth-three circle campaign also measures nontrivial hidden-feature motion,
so the agreement is not explained by a frozen-feature limit.  Experiments with
GELU and sigmoid show that the construction is not tanh-specific, but accuracy
is task dependent: the reported GELU hard-outlier $P=3$ error is $0.02485$,
whereas harder GELU/sigmoid cases are near $0.26$ and $0.24$.  ReLU and SELU
runs are useful empirical probes but fall outside the smooth ODE theorem at
their kinks; selected ReLU flows can even be nonunique at exceptional initial
states.

\begin{figure}[t]
 \centering
 \begin{subfigure}[t]{\textwidth}
  \centering
  \includegraphics[width=\textwidth]{figures/mnist_scatter.pdf}
  \caption{Closure versus dense predictions on 1984 held-out images.}
 \end{subfigure}
 \vspace{0.6em}
 \begin{subfigure}[t]{0.75\textwidth}
  \centering
  \includegraphics[width=\textwidth]{figures/mnist_rms.pdf}
  \caption{RMS discrepancy across several matched-loss endpoints.}
 \end{subfigure}
 \caption{MNIST 3-versus-8, 100 training images, width $4096$.  At training MSE
 $0.001$, $P=1,2,3$ have held-out RMS discrepancies
 $0.00324756,0.00108443,0.00119922$.  All declared numerical-refinement gates
 passed; the slight $P=3$ worsening shows that autonomous feedback need not make
 realized error monotone in $P$.}
 \label{fig:mnist}
\end{figure}
\FloatBarrier

For this MNIST panel the learned rank bounds are $100,200,300$.  The reported
learned-increment storage is 6.25, 12.5 and 18.75 MiB, versus 128 MiB for the
unrestricted learned internal correction.  The closure also retains the
128 MiB initialized matrix, so its minimum total state is not smaller than a
dense implementation storing only current weights.  The measured adaptive
closure integration was also slower in this panel.  The empirical achievement
is fidelity with a compressed moving update, not an end-to-end speedup claim.

\section{Limitations and next questions}

The main positive result is a fixed-order family of autonomous systems whose
state does not grow with elapsed training time and whose physical trajectories
converge to dense gradient flow as $P\to\infty$ on every prescribed finite
horizon.  Several boundaries matter.

\begin{itemize}
\item The initialized random matrices and their transpose actions are retained.
Compressing that source while preserving matrix reuse is still open.
\item Constants are not uniform in width, depth, sample count or horizon.  A
direct fixed-$P$ population theorem and a CLT-style finite-width error theory
remain future work.
\item Neither theorem gives one order valid for all $t\ge0$.  Uniform all-time
tracking would require bounded accumulated activity and nonlinear stability
with finite total amplification.
\item Smooth activations such as tanh, sigmoid, arctangent, GELU, softplus,
SiLU, sine and polynomials fall within the stated local smoothness classes.
Literal ReLU and standard SELU do not.  Their discontinuous selected backward
fields require nonsmooth dynamics and a clock that charges jump variation.
\item A finite scalar aggregate truncation is a second approximation, beyond
history truncation.  Differentiating aggregates gives a principled hierarchy,
but an economical width-uniform scalar closure is not established here.
\item Matching the dense predictor is different from predicting the target
well.  Generalization requires an additional argument about the function
selected by training.
\end{itemize}

The response-memory construction nevertheless gives a concrete answer to the
paper's question.  Exact dynamics can hide a growing history inside a formal
population or operator description.  Approximate dynamics can replace that
history by finitely many evolving coordinates per neuron, retain nonlinear
feature feedback, and admit a genuine order-versus-trajectory-error theorem.
That smaller state is a new lens on the learned part of deep training, even
before the remaining random operator is compressed.

\bibliographystyle{plainnat}
\bibliography{references}

\end{document}
